% GENERATED_BY: oc_core_monograph_translation_draft_pipeline_v1.py
% AI_ASSISTED_TRANSLATION_DRAFT: TRUE
% THEOREM_REVIEW_STATUS: PENDING
% THEOREM_REVIEW_MODE: FORMAL_AI_GATE__CODEX_CLI_CHATGPT
% THEOREM_REVIEW_REVIEWER_ID:
% THEOREM_REVIEW_AUTHORITY_CLASS:
% THEOREM_REVIEWED_AT:
% SOURCE_LANGUAGE: EN
% TARGET_LANGUAGE: RU
% SOURCE_EN_REF: content/m_spaces/m9.tex
% ================================================================
% ==== FILE: content/m_spaces/m9.tex
% ================================================================

% ==============================
%  Ontology of Continua — Core
%  Meta-Space: M9
%  FULL MODULE — FINAL VERSION
% ==============================

\subsubsection{\texorpdfstring{$M_9$}{M_9} Обзор}
\label{sec:m9-overview}

$M_9$ — мета-пространство, обеспечивающее существование
теоретических континуумов ($K_9$).
В то время как $M_8$ задаёт материально-символическую и институциональную
подоснову цивилизаций, $M_9$ вводит структурные условия для
возникновения, сосуществования и преобразования
научных, математических и концептуальных моделей.

В $M_9$ базовыми единицами являются не агенты, инфраструктуры и
институты, а \emph{модели} — структурированные объекты,
имеющие внутренние логики, ограничения и правила преобразований.

Следовательно, $M_9$ — минимальное пространство, в котором
\emph{теория} становится допустимым континуумом.

% ---------------------------------------------------------------
\subsubsection*{1. Допустимое пространство конфигураций \texorpdfstring{$\Omega(M_9)$}{\Omega(M_9)}}

\[
\Omega(M_9)=
\left\{
(\mathcal{M},\ \mathrm{Mor},\ A_{\mathrm{model}},\
\Theta_{\mathrm{coh}},\ C_{\mathrm{theory}},\ \Sigma_9)
\mid \text{условия когерентности выполняются}
\right\}.
\]

Компоненты:

\begin{itemize}
    \item $\mathcal{M}$ — множество допустимых моделей
        (математических, физических, химических, биологических, когнитивных, социальных),
    \item $\mathrm{Mor}$ — морфизмы между моделями
        (интерпретации, редукции, эквивалентности, функторы),
    \item $A_{\mathrm{model}}$ — оси теоретической вариативности,
    \item $\Theta_{\mathrm{coh}}$ — пороги когерентности,
    \item $C_{\mathrm{theory}}$ — циклы теоретического уточнения,
    \item $\Sigma_9$ — глобальные структурные ограничения пространства теории.
\end{itemize}

Допустимость требует:

\begin{enumerate}
    \item существования устойчивых моделей: $\mathcal{M}\neq\emptyset$,
    \item существования содержательных морфизмов между моделями,
    \item нетривиальных теоретических осей,
    \item когерентности: $d_{\mathrm{coh}}(M_i, M_j)<\Theta_{\mathrm{coh}}$ для допустимых пар,
    \item возможности теоретической эволюции под действием оператора $E$.
\end{enumerate}

% ---------------------------------------------------------------
\subsubsection*{2. Оси \texorpdfstring{$A(M_9)$}{A(M_9)}}

\[
A(M_9)=
\{
A_{\mathrm{syntax}},\
A_{\mathrm{semantics}},\
A_{\mathrm{abstraction}},\
A_{\mathrm{idealization}},\
A_{\mathrm{approximation}},\
A_{\mathrm{domain}},\
A_{\mathrm{logic}}
\}.
\]

Определения:

\begin{itemize}
    \item $A_{\mathrm{syntax}}$ — формальные языки, символические структуры,
    \item $A_{\mathrm{semantics}}$ — пространства интерпретации,
    \item $A_{\mathrm{abstraction}}$ — степени обобщённости,
    \item $A_{\mathrm{idealization}}$ — допустимые упрощения,
    \item $A_{\mathrm{approximation}}$ — ошибки и структуры сходимости,
    \item $A_{\mathrm{domain}}$ — оси, специфичные для предметной области (физика, химия, биология и т.д.),
    \item $A_{\mathrm{logic}}$ — логические правила, системы вывода.
\end{itemize}

Необходимое условие для модельного существования:

\[
\dim(A_{\mathrm{syntax}})\ge 1,
\quad
\dim(A_{\mathrm{semantics}})\ge 1.
\]

% ---------------------------------------------------------------
\subsubsection*{3. Потенциалы \texorpdfstring{$P(M_9)$}{P(M_9)}}

\[
P(M_9)=
(F_{\mathrm{expressivity}},\
F_{\mathrm{consistency}},\
F_{\mathrm{coherence}},\
F_{\mathrm{generalization}},\
F_{\mathrm{predictive}},\
F_{\mathrm{compression}}).
\]

Интерпретация:

\begin{itemize}
    \item $F_{\mathrm{expressivity}}$ — потенциал, обеспечивающий богатые теоретические описания,
    \item $F_{\mathrm{consistency}}$ — потенциал, обеспечивающий логическую устойчивость,
    \item $F_{\mathrm{coherence}}$ — потенциал, сохраняющий взаимную интерпретируемость моделей,
    \item $F_{\mathrm{generalization}}$ — способность к абстракции и переносу между областями,
    \item $F_{\mathrm{predictive}}$ — точность и эмпирическая достаточность,
    \item $F_{\mathrm{compression}}$ — потенциал для минимально достаточных представлений.
\end{itemize}

Условия существования $K_9$:

\[
F_{\mathrm{coherence}}>\Theta_{\mathrm{fragmentation}},
\qquad
F_{\mathrm{consistency}}>\Theta_{\mathrm{paradox}}.
\]

% ---------------------------------------------------------------
\subsubsection*{4. Пороги \texorpdfstring{$\Theta(M_9)$}{\Theta(M_9)}}

Ключевые пороги когерентности:

\begin{itemize}
    \item $\Theta_{\mathrm{coh}}$ — максимальное допустимое расхождение для сосуществования моделей,
    \item $\Theta_{\mathrm{consistency}}$ — порог внутренней логической устойчивости,
    \item $\Theta_{\mathrm{predictive}}$ — минимальная предикативная точность,
    \item $\Theta_{\mathrm{paradox}}$ — граница, где несогласованность приводит к коллапсу континуума,
    \item $\Theta_{\mathrm{fragmentation}}$ — порог теоретической дезинтеграции.
\end{itemize}

Превышение $\Theta_{\mathrm{paradox}}$ приводит к схлопыванию $\Omega(K_9)$,
аналогично институциональному коллапсу в $K_8$.

% ---------------------------------------------------------------
\subsubsection*{5. Потоки \texorpdfstring{$J(M_9)$}{J(M_9)}}

\[
J(M_9)=
(J_{\mathrm{inference}},\
J_{\mathrm{interpretation}},\
J_{\mathrm{translation}},\
J_{\mathrm{abstraction}},\
J_{\mathrm{specialization}},\
\nabla_i T_9).
\]

Интерпретация:

\begin{itemize}
    \item $J_{\mathrm{inference}}$ — поток дедукций и доказательств,
    \item $J_{\mathrm{interpretation}}$ — переинтерпретация между моделями,
    \item $J_{\mathrm{translation}}$ — перевод между формальными языками,
    \item $J_{\mathrm{abstraction}}$ — поток к более высокой степени общности,
    \item $J_{\mathrm{specialization}}$ — предметно-специфические выводы,
    \item $\nabla_i T_9$ — градиенты теоретического натяжения.
\end{itemize}

Условие устойчивости:

\[
\oint J_{\mathrm{inference}}\,dt\approx 0
\quad\Rightarrow\quad
\text{модель остается логически устойчивой на циклах}.
\]

% ---------------------------------------------------------------
\subsubsection*{6. Циклы \texorpdfstring{$C(M_9)$}{C(M_9)}}

Фундаментальные циклы:

\begin{itemize}
    \item $C_{\mathrm{theory}}$
        — гипотеза $\to$ вывод $\to$ оценка $\to$ пересмотр,
    \item $C_{\mathrm{abstraction}}$
        — конкретное $\to$ абстрактное $\to$ обобщённое $\to$ инстанцированное,
    \item $C_{\mathrm{interpretation}}$
        — модель $\to$ отображение $\to$ переинтерпретация,
    \item $C_{\mathrm{coherence}}$
        — локальные соответствия $\to$ глобальная согласованность $\to$ циклы подгонки,
    \item $C_{\mathrm{unification}}$
        — модели области $\to$ междоменная структура $\to$ единая рамка.
\end{itemize}

Живой теоретический континуум требует:

\[
\oint_{C_{\mathrm{coherence}}} dA \approx 0.
\]

% ---------------------------------------------------------------
\subsubsection*{7. Граница \texorpdfstring{$\partial\Omega(M_9)$}{\partial\Omega(M_9)}}

Условия приближения к границе:

\begin{itemize}
    \item разрушение замыкания вывода,
    \item потеря когерентности: $d_{\mathrm{coh}}>\Theta_{\mathrm{coh}}$,
    \item лавинообразное появление парадоксов,
    \item расходимость формальных языков за пределы взаимной интерпретируемости,
    \item потеря соответствия модель–реальность (предиктивный коллапс),
    \item разгон теоретического натяжения $T_9$.
\end{itemize}

Пересечение $\partial\Omega(M_9)$ приводит к свёртыванию теории в фрагментированные структуры, допустимые в $M_8$.

% ---------------------------------------------------------------
\subsubsection*{8. Связь с \texorpdfstring{$M_8$}{M_8} и $M_{10}$}

\textbf{От $M_8$ к $M_9$:}

Необходимые условия:

\[
A_{\mathrm{syntax}}\neq\emptyset,
\quad
A_{\mathrm{semantics}}\neq\emptyset,
\quad
\Theta_{\mathrm{coh}}\ \text{finite}.
\]

Символические и научные структуры $M_8$ поднимаются в
пространства моделей $M_9$ оператором $\Psi$.

\textbf{К $M_{10}$ (мета-теоретическое пространство):}

$M_{10}$ вводит:

\begin{itemize}
    \item категории категорий моделей,
    \item функториальные динамики,
    \item пороги мета-кохерентности,
    \item циклы самописания,
    \item пространство, где сами метатеории образуют континуумы.
\end{itemize}

Таким образом, $M_9$ задаёт упорядоченное модельное пространство, которое становится подосновой для $K_{10}$.
